\begin{helper}[Identification of corresponding PE witnesses]
Let two PE constructions be made on the same ordered tree \(T\), with
maximum-defective choices
\((v,u,f,\delta)\) and \((v^*,u^*,f^*,\delta^*)\), and extra colors
\(\gamma\) and \(\gamma^*\).  Suppose their singleton connecting-edge
sets and unordered stored-color pairs agree.  Suppose also that their
post-switch missing-color sets agree at the two support vertices, that each
delta is missing there, and that each gamma is present there.  Then
\[
 f^*=f,\qquad v^*=v,\qquad u^*=u,\qquad
 \delta^*=\delta,\qquad\gamma^*=\gamma .
\]
\end{helper}
\begin{proof}
Equality of the singleton edge sets gives \(f^*=f\) as an exact edge
identity.  The endpoint equation in each
\noderef{MaximumDefectiveChoice} writes the endpoints of this edge as
\(\{v,u\}\) and \(\{v^*,u^*\}\).  Both support vertices belong to
the common tree, whereas both other endpoints lie outside it.  The crossed
orientation would put an outside endpoint in the tree, so the endpoint
equality has the unique orientation \(v^*=v\) and \(u^*=u\).

The unordered pair equality says
\(\{\delta,\gamma\}=\{\delta^*,\gamma^*\}\).
Under the identified support vertex, the common post-switch missing set
contains \(\delta\) and \(\delta^*\), but contains neither
\(\gamma\) nor \(\gamma^*\).  If \(\delta^*=\gamma\), the same
color would be both missing and present, a contradiction.  Hence the
two-element-set equality forces \(\delta^*=\delta\).  Substituting this
equality back into the pair equality, the alternative
\(\gamma^*=\delta\) is ruled out in the same way, leaving
\(\gamma^*=\gamma\).  This proves all five identities.
\end{proof}

